\chapter{监测资源优化策略与研究展望}\label{ch:policy}

\section{国家公共卫生监测体系的资源配置策略矩阵}\label{sec:policy_matrix}

基于本文建立的理论推断框架、数值模拟发现与全美真实流感及欧洲 COVID-19 实证检验，本节针对国家公共卫生疾控部门在传染病监测预警网络建设中的资源配置，提出一套可落地的**“两阶段差异化监测资源优化配置策略矩阵”**。

以往的监测资源配置往往采取“一刀切”的机械扩容策略，在疫情爆发时一味要求所有定点医院加班加点扩大采样量，在疫情平缓期则疏于维持基础检测网络。本文的研究证明，这种策略无论在统计效率上还是经济成本上都是低效的。正确的资源配置逻辑必须**依循当前系统规模与内生特征平衡点 $m_\times$ 的相对位置动态调整**。

\subsection{流行初期与低发期 ($m < m_\times$) ——「全力扩量，迅速越峰」}
在传染病暴发早期、散发期或换季低发期，医疗系统实际捕获的病例数稀疏，系统处于 $m < m_\times$ 的**微观抽样噪声主导体制**：
\begin{itemize}
    \item **核心矛盾**：此时微观泊松相遇离散与超级传播重尾波动占总推断方差的 50\% 以上。若哨点过少，少量超级传播者的出现或缺席会导致再生数估计发生剧烈抖动，常规正态检验频繁出现假警报（如第四章揭示的 8.1\% 虚警率）；
    \item **配置策略**：此时继续扩大病例捕获规模具有极高的统计边际收益（投入产出比极高）。卫生部门应集中财政预算扩建定点哨点网络、降低就诊检测门槛、提高轻症筛查覆盖率 $\rho$ 与地理人群基数 $n$，推动系统尽快跨越特征平衡点 $m_\times$；
    \item **预警指标优化**：在进行疫情爆发早期检测时，依据目标病原体的超离散度参数 $k$，精确测算达到 80\% 功效所需的充分样本门槛 $m_{80\%}$，作为低发期哨点捕获能力的刚性考核指标。
\end{itemize}

\subsection{流行达峰与饱和期 ($m > m_\times$) ——「战略提质，协同平滑」}
当疫情全面暴发、每周报告病例数显著超过特征平衡点（$m \gg m_\times$）时，系统进入**宏观环境随机性饱和体制**：
\begin{itemize}
    \item **核心矛盾**：大数定律的边际降差速率急剧放缓。总方差由全社会共同面临的外生环境冲击（不可约环境底板 $v_R$）以及疫情转折导致的基准漂移（占达峰期误差的 88.9\%）完全主导。此时继续强行要求医院加倍采样上报，在统计上几乎无法带来实质性的预测精度提升；
    \item **配置策略**：监测资源必须坚决由“扩量”转向“提质”，具体实施三项战略转移：
    \begin{enumerate}[label=(\arabic*)]
        \item **引入多模型前瞻集成系统**：如第六章 FluSight 实证所示，官方自适应集成预测可将达峰期均值偏差由 88.9\% 剧烈压缩至 15.5\%，这是单模型狂堆样本量根本无法做到的；
        \item **采集宏观环境协变量**：将预算转向监测人群宏观流动大数据、气象降温湿变指数与学校企事业单位出勤率，通过在状态方程中显式控制环境协变量，直接压缩不可约环境方差 $v_R$ 本身；
        \item **空间多中心弱相关汇聚**：如第五章实证所示，跨区域大数加总能够利用区域间环境扰动的弱相关性（$\bar{c} < 1$），在空间维度实现类似时间平均的降差收益。
    \end{enumerate}
\end{itemize}

\section{研究局限性与方法学边界}\label{sec:limitations}

为了恪守严谨求实的现代学术规范，本节坦诚剖析本研究存在的方法学局限与适用边界：

\subsection{不可约环境方差 $v_R$ 始终为拟合量而非物理直接测量量}
在第三章与第五章中，理论上的环境方差底板 $v_R$ 在实证中对应 FGLS 仿射拟合模型的截距项 $\hat{a}$。必须承认，在真实的医院住院监测数据中，截距项 $\hat{a}$ 是一个唯象的复合量：
它既包含了天气、人群行为等纯外生宏观环境方差，也不可避免地混合了由于滞后卷积权重失配、报告周期截断带来的模型设定偏差。由于缺乏独立的对照实验，实证中的经验平衡点 $\hat{b}/\hat{a}$ 应当理解为具体监测系统在当前技术条件下的**量级参照物**，而非类似于物理常数的绝对点位识别。

\subsection{观测覆盖率变动与事件级二项抽稀的微观机制差异}
第五章在评估 2024 年自愿填报政策冲击时，实测数据显示医院覆盖率骤降至 54.6\% 引发了分布右尾的极端扩张（偏度激增至 14.98），而四分位距与中位绝对离差基本平稳。这一现象表明，在现实制度变迁中，医院填报率的下滑往往是“整所医院直接退出直报系统”（整群非随机缺失），而非理想模型中所设想的“全社会每个患者独立抛硬币决定是否上报”（均匀事件级二项抽稀）。虽然二者在放大观测方差的大方向上定性一致，但在更高阶矩的分布形态上存在微观机制差异。

\subsection{时间序列卷积中的单向因果性与内生性}
在更新过程时间序列预测中，构建的自变量为历史滞后有效暴露规模 $D_t = \sum w_s C_{t+1-s}$。虽然在时间维度上严格保证了“历史发生于当期之前”的事前信息集顺序，排除了前瞻偏误；但传染病的时间序列具有强烈的内生自回归性，严格的因果识别仍需借助更丰富的外生工具变量或断点回归设计。

\section{未来研究展望：双轨前瞻主动追踪队列的解耦设计}\label{sec:future_work}

针对上述局限性，本文在数理流行病学与监测设计方法学上提出一个极具操作可行性的未来拓展方向——**“双轨主动追踪队列识别设计”（Dual-track Prospective Cohort Design）**：

目前全球绝大多数疾控系统均采用单一的“被动医院住院就诊直报网络”。医院数据天生混合了就医意愿、医保报销政策与填报时滞，导致外生环境方差与系统观测偏差永远纠缠在一起。
未来的国家公共卫生监测体系可以在保留大型医院被动住院监测网络（第一轨）的同时，在代表性社区布设**固定人群规模的前瞻性主动追踪队列（第二轨）**：
\begin{itemize}
    \item 在不同地理气候带抽取具有人口统计学代表性的固定家庭队列；
    \item 无论家庭成员是否有症状，定期进行抗原筛查与多病原体核酸检测，排除就医意愿干扰，直接测得无偏差的真实社区传播感染率；
    \item 通过联立两轨监测数据，利用计量经济学联立方程组或状态空间反卷积算法，即可将**纯外生宏观环境方差 $v_R$** 与**医院监测系统的代理平滑偏差**彻底解耦识别！
\end{itemize}

\section{本章小结与全文总结}\label{sec:conclusion}

本学位论文围绕传染病监测中“数据规模崇拜”这一现实科学悖论，从随机分支过程、全方差分解与抽样运筹理论出发，完成了系统的数理推导、仿真检验与大数据实证：
\begin{enumerate}[label=(\arabic*)]
    \item **破除教条**：严格确立了“推断目标决定规模收益”的核心规律。证实当期瞬态点估计（目标 A）规模收益无界（$O(1/m)$），而跨期基准估计（目标 B）必然受制于不可约环境方差底板 $v_R$ 与特征规模平衡点 $m_\times$；
    \item **运筹对偶**：证明了提高报告率与扩大基数的非对称性，揭示了样本预算恒常性，导出了显式经费约束下最优单期规模闭式解及其与 Cochran 两阶段抽样模型的数学同构；
    \item **实证反转**：利用全美六年 16,218 条真实住院数据（10,303 活跃周样本），证实规模反比项高度显著为正（$\hat{b} = 2.24$）；揭示达峰期 88.9\% 误差来源于基准漂移，有力推翻了“达峰期样本量失效”的传统误读；
    \item **外验证实**：通过 CDC FluSight 官方集成预测证明多模型集成能将达峰偏差剧烈压缩至 15.5\%，欧洲 COVID-19 数据证实了规律的跨病种普适性与尺度依赖性。
\end{enumerate}

综上所述，本文为国家公共卫生监测体系在不同流行阶段由“盲目扩量”迈向“科学提质”提供了完备的数理统计基础与量化决策工具。


\section{面向疾控机构的监测资源优化配置量化操作手册}\label{sec:actionable_manual}

为了将前文较为抽象的数学推导与计量参数切实转化为公共卫生各级疾控中心（CDC）在日常与应急状态下可直接执行的决策工具，本节将理论框架细化为四个标准化步骤的操作工作流。

\subsection{第一步：监测系统基线参数测定与建档}
在每年呼吸道流行季来临前（如每年 9 月初），辖区疾控监测部门应首先对本地区的历史监测系统进行基线参数标定：
\begin{enumerate}[label=(\arabic*)]
    \item **确定目标病原体生物学离散参数 $k$**：依据第二章表~\ref{tab:superspreading_meta}，对季节性流感设定基准先验 $k = 0.35$；对新型冠状病毒设定 $k = 0.20$；
    \item **评估哨点医院网络复合报告率 $\rho$**：结合本辖区常住人口、医保就诊率与哨点定点医院覆盖比例，核算系统报告率 $\rho$；
    \item **测算本辖区特征规模平衡点 $m_\times$ 与预警门槛 $m_{80\%}$**：
    利用公式 $m_\times = \frac{\bar{R} + \rho(\bar{R}^2+v_R)/k}{v_R}$ 与命题 5，计算出本地区的特征平衡点 $m_\times$（周度流感住院量级一般为 20--30 例/周），以及增长检测 80\% 功效门槛 $m_{80\%}$。
\end{enumerate}

\subsection{第二步：暴露规模在线滑动卷积监控}
在疫情监测运转期间，信息系统应自动以周为步长，对最近 3 周的报告住院数执行代际卷积平滑：
\begin{equation}
D_t = 0.65 C_t + 0.25 C_{t-1} + 0.10 C_{t-2}.
\end{equation}
系统将当周计算出的有效规模 $D_t$ 与预设的特征平衡点 $m_\times$ 进行实时布尔比对，触发两阶段资源调度开关。

\subsection{第三步：两阶段资源调度自动切换规则}
系统依据 $D_t$ 与 $m_\times$ 的相对位置，严格实施以下调度切换：
\begin{itemize}
    \item **若 $D_t \le m_\times$（微观噪声主导区，执行“扩量预案”）**：
    \begin{enumerate}[label=\alph*.]
        \item 财政预算重点投向扩建发热门诊检测通道、增加基层社区采样频次；
        \item 确保每周报告数尽快达到 $m_{80\%}$，防止小计数导致的虚警与漏报；
        \item 暂缓投入昂贵的多模型超级计算预测系统，因为在小样本下任何复杂算法都会被抽样噪声淹没。
    \end{enumerate}
    \item **若 $D_t > m_\times$（宏观环境饱和区，执行“提质预案”）**：
    \begin{enumerate}[label=\alph*.]
        \item 立即对各哨点医院发出采样减负指令，停止盲目扩充同质化的简单病例计数报告；
        \item 启动前瞻多模型自适应集成预报平台，利用模型多样性主动校准达峰期的均值漂移（消解 88.9\% 偏差）；
        \item 紧急调取交通卡口人口流动、手机信令、气象降温及抗原自测上报等外部环境协变量，显式输入预测模型以平滑环境底板 $v_R$；
        \item 启动与相邻省市疾控中心的跨区域空间数据汇聚机制，利用空间弱相关汇聚效应拉低宏观预测方差。
    \end{enumerate}
\end{itemize}

上述量化工作流将本文复杂的理论命题完全转化为了具有明晰逻辑判断、参数输入与动作指令的现代疾控决策管理闭环。
